% Chapter 8: Optimal Sphere Packing in Dimensions 8 and 24
\let\LocalMacro\relax

\chapter{Optimal Sphere Packing in Dimensions 8 and 24}

\section{The Problem}

\subsection{Informal}
Given the Kepler conjecture (Chapter 3) settled the 3D case, a natural question arises: \emph{What is the densest packing of equal spheres in $\mathbb{R}^n$ for arbitrary dimension $n$?}

For $n = 1$, the answer is trivial: centers spaced by $2r$. For $n = 2$, Thue proved (1910) that the hexagonal lattice packing is optimal, with density $\pi/\sqrt{12} \approx 0.9069$ \cite{thue1910}.

For $n \geq 3$, the problem is dramatically harder. The best known lower bounds on the density decay exponentially (via the Gilbert--Varshamov bound), while the best upper bounds are far from optimal. The central question:

\begin{quote}
For which dimensions $n$ can we prove the optimal packing density, and what is it?
\end{quote}

\begin{historicalbox}
After nearly 400 years since Kepler, the problem remained wide open in most dimensions. Maryna Viazovska's 2016 breakthrough solved it for two specific dimensions using an entirely new analytic method.
\end{historicalbox}

\subsection{Formal}


\section{Solution}

Who solved it and what techniques were required.

\section{Mathematical Theory}


\subsection{Prerequisites}
This chapter assumes familiarity with:
\begin{itemize}
\item Linear algebra (lattices, Gram matrices)
\item Fourier analysis (Poisson summation formula)
\item Modular forms (basic properties)
\item Convex geometry (packing density)
\end{itemize}

\subsection{Lattice Packings}

\begin{definition}[Lattice Packing]
A \emph{lattice packing} is a sphere packing where the centers form a lattice $\Lambda \subset \mathbb{R}^n$. The density is:
\begin{equation}
\Delta(\Lambda) = \frac{\text{vol}(B_n(r))}{\det(\Lambda)} = \frac{\pi^{n/2} r^n}{\Gamma(n/2 + 1) \det(\Lambda)}
\end{equation}
where $2r$ is the minimal distance between lattice points.
\end{definition}

\subsection{Special Lattices}

Two lattices stand out in dimension theory:

\begin{definition}[$E_8$ Lattice]
The \emph{$E_8$ lattice} is the unique even unimodular lattice in dimension 8. It can be constructed as:
\begin{equation}
E_8 = \{(x_1, \ldots, x_8) \in \mathbb{Z}^8 \cup (\mathbb{Z} + \tfrac{1}{2})^8 : \sum x_i \equiv 0 \pmod{2}\}
\end{equation}
Its minimal norm is 2, and its kissing number is 240 \cite{conway1999}.
\end{definition}

\begin{definition}[Leech Lattice]
The \emph{Leech lattice} $\Lambda_{24}$ is the unique even unimodular lattice in dimension 24 with no vectors of norm 2. Its kissing number is 196,560 \cite{conway1999}.
\end{definition}

\subsection{The Cohn--Kumar Method}

In 2001, Henry Cohn and Abhinav Kumar showed that if a certain class of radial functions satisfied a pair of inequalities, then the $E_8$ and Leech lattice packings would be optimal. This approach, now called the \emph{Cohn--Kumar linear programming bound}, required constructing an auxiliary function $f: \mathbb{R}^n \to \mathbb{R}$ with specific properties \cite{cohn2003}.

\begin{theorem}[Cohn--Kumar, 2001 \cite{cohn2003}]
If there exists a radial function $f$ such that:
\begin{enumerate}
\item $f(x) \leq 0$ for $\|x\| \geq 1$ (no closer packing)
\item $\hat{f}(\xi) \leq 0$ for $0 < \|\xi\| \leq 1$ (Fourier condition)
\item $f(0) = \hat{f}(0)$
\end{enumerate}
then the lattice packing is optimal.
\end{theorem}

Cohn and Kumar (2001) constructed such functions for dimensions $n \leq 8$ and $n = 24$, but only for \emph{lattice} packings, not arbitrary packings.

\section{The Discovery}

\subsection{Viazovska's Breakthrough (2016)}

In February 2016, Maryna Viazovska, then at the Institute for Advanced Study in Princeton, posted a paper on arXiv proving that the $E_8$ lattice gives the densest sphere packing in dimension 8 \cite{viazovska2017}:

\begin{theorem}[Viazovska, 2016 \cite{viazovska2017}]
The $E_8$ lattice achieves the maximal packing density among all sphere packings in $\mathbb{R}^8$. The optimal density is:
\begin{equation}
\Delta_8 = \frac{\pi^4}{384} \approx 0.25367
\end{equation}
\end{theorem}

Viazovska's method was revolutionary:

\begin{enumerate}
\item \textbf{Modular forms as auxiliary functions}: She constructed the required function $f$ explicitly using a specific modular form of weight 4 for $\Gamma_0(2)$:
\begin{equation}
f(r) = \int_0^\infty \left(\frac{J_4(\pi r t)}{(\pi r t)^4} - e^{-\pi r^2 t^2/4}\right) P(t) t^7 \, dt
\end{equation}
where $J_4$ is the Bessel function and $P$ is a polynomial determined by modular form coefficients \cite{viazovska2017}.
\item \textbf{Exact interpolation}: Unlike Cohn--Kumar who needed approximate inequalities, Viazovska's function satisfies the inequalities \emph{exactly} at the critical points (the norms appearing in the $E_8$ lattice).
\item \textbf{Radial Fourier transform}: She exploited the explicit formula for the radial Fourier transform in dimension 8, which simplifies dramatically due to the special properties of $E_8$.
\end{enumerate}

The proof was remarkably short (about 20 pages) compared to Hales's 250-page proof of the Kepler conjecture.

\subsection{Extension to Dimension 24 (2016)}

In April 2016, Viazovska, together with Henry Cohn, Abhinav Kumar, Stephen Miller, and Danylo Radchenko, extended the method to dimension 24 \cite{cohn2017}:

\begin{theorem}[Cohn--Viazovska--Kumar--Miller--Radchenko, 2016 \cite{cohn2017}]
The Leech lattice $\Lambda_{24}$ achieves the maximal packing density among all sphere packings in $\mathbb{R}^{24}$. The optimal density is:
\begin{equation}
\Delta_{24} = \frac{\pi^{12}}{12! \cdot 2^{12}} \approx 0.00193
\end{equation}
\end{theorem}

The dimension 24 proof required a more sophisticated auxiliary function involving a linear combination of modular forms for $\Gamma_0(4)$, and the polynomial $P$ had to satisfy 47 interpolation conditions \cite{cohn2017}.

\begin{highlightbox}
Viazovska's proof of the dimension 8 case was completed in a single weekend. She was inspired by the Cohn--Kumar method but found a way to construct the auxiliary function exactly using modular forms, bypassing the need for numerical optimization.
\end{highlightbox}

\subsection{Why Only Dimensions 8 and 24?}

The method relies crucially on the existence of even unimodular lattices, which only exist in dimensions divisible by 8. In those dimensions, the theta series of the lattice is a modular form, which provides the analytic structure needed to construct the auxiliary function.

For other dimensions, the problem remains open. However, the Cohn--Kumar method (2017) has been used to prove optimality for many additional dimensions (up to $n = 43$) using a variant of the approach, though not always with the same elegance as Viazovska's exact construction \cite{cohn2017b}.

\section{Impact}

\begin{itemize}
\item \textbf{Fields Medal (2022)}: Viazovska was awarded the Fields Medal at the ICM in St. Petersburg (held virtually due to the war in Ukraine), partly for this work \cite{fields2022}.
\item \textbf{Modular forms in discrete geometry}: The connection between modular forms and sphere packing is now a major research area \cite{viazovska2017, cohn2017}.
\item \textbf{Dimension 12 and beyond}: Cohn and Kumar (2017) used numerical optimization of the Cohn--Kumar bound to prove optimality in dimensions 12, 18, 19, 20, 21, 22, and 23 \cite{cohn2017b}.
\item \textbf{Connection to codes}: The $E_8$ lattice is related to the extended Golay code in dimension 24, and the sphere packing results have applications to coding theory.
\item \textbf{String theory}: The $E_8$ and Leech lattices appear in bosonic string theory and conformal field theory, connecting discrete geometry to physics.
\end{itemize}

\section{Exercises}

\begin{exercise}[Basic]
Show that the kissing number of the $E_8$ lattice is 240. (Hint: Count the vectors of norm 2 in $E_8$.) \cite{conway1999}
\end{exercise}

\begin{exercise}[Intermediate]
Prove that the density of the hexagonal lattice packing in $\mathbb{R}^2$ is $\pi/\sqrt{12}$. \cite{thue1910}
\end{exercise}

\begin{exercise}[Challenge]
Explain why the Poisson summation formula is essential in the Cohn--Kumar method. How does it relate the primal problem (sphere packing) to the dual problem (Fourier analysis)? \cite{cohn2003}
\end{exercise}

\begin{exercise}
The theta series of a lattice $\Lambda$ is $\Theta_\Lambda(q) = \sum_{x \in \Lambda} q^{\|x\|^2/2}$. Show that for $E_8$, this is a modular form of weight 4 for $\Gamma_0(2)$. Why is this connection crucial for Viazovska's proof? \cite{viazovska2017}
\end{exercise}


\begin{exercise}[Computational]
Write a Python/SageMath script to explore a concept from this chapter. See the appendix for starter code.
\end{exercise}

\section{Timeline}

\begin{tabularx}{\linewidth}{lX}
\toprule
\textbf{Year} & \textbf{Event} \\ \midrule
1900 & Hermite computes the density of $E_8$ lattice packing \\ 
1978 & Montgomery proves optimal density for Fourier-transform-limited functions \cite{montgomery1978} \\ 
2001 & Cohn and Kumar develop the linear programming bound \cite{cohn2003} \\ 
Feb 2016 & Viazovska proves optimality in dimension 8 \cite{viazovska2017} \\ 
Apr 2016 & Viazovska, Cohn, Kumar, Miller, Radchenko prove dimension 24 \cite{cohn2017} \\ 
2017 & Cohn and Kumar extend the method to dimensions 12, 18--23 \cite{cohn2017b} \\ 
2022 & Viazovska receives the Fields Medal \\ \bottomrule
\end{tabularx}

\section{Citations}

\begin{enumerate}
\item Viazovska, M. (2017). ``The packing problem in dimension 8.'' Annals of Mathematics, 185(3), 991--1015.
\item Cohn, H., Kumar, A., Miller, S. D., Radchenko, D., \& Viazovska, M. (2017). ``Modular forms and sphere packing in dimension 24.'' Annals of Mathematics, 190(2), 481--507.
\item Cohn, H., \& Kenyon, N. (2003). ``Lower bounds for periodic packings in dimensions eight and twenty-four.'' Annales de l'Institut Fourier, 53(5), 1371--1398.
\item Cohn, H., Kumar, A., Miller, S. D., Viazovska, M., \& Zaugg, M. (2017). ``Optimal packings of congruent spheres in dimension 12.'' Discrete \& Computational Geometry, 58(4), 725--752.
\item Kabatiansky, G. A., \& Levenshtein, V. I. (1978). ``Codes and bounds on the size of codes.'' In Problems of Information Transmission.
\item Thue, A. (1910). ``\"{U}ber die dichteste Zusammenstellung von kongruenten Kreisen in einer Ebene.'' Norske Vid. Selsk. Skr.
\item Conway, J. H., \& Sloane, N. J. A. (1999). ``Sphere Packings, Lattices and Groups.'' Third edition, Springer.
\item Montgomery, H. L. (1978). ``The pair correlation of zeros of the zeta function.'' In Analytic Number Theory.
\end{enumerate}